% GENERATED FILE. Edit canonical lessons, then run npm run generate:books.
% canonical-source-hash: fnv1a64:1bf769e3ce292963
% canonical-lessons: FR-C338-douter, FR-C338-sceptique, FR-C338-nier, FR-C338-absurde, FR-C338-discutable

\chapter{Je n'y crois pas}
\label{ch:fr-je-n-y-crois-pas}
\hlchaptermodality{338}

\begin{chapteropening}
\textbf{By the end of this chapter:} \emph{I can say I doubt something, that I am sceptical, deny a claim and say an idea is absurd or debatable.}

Complete the last lesson of chapter 338: I can say I doubt something, that I am sceptical, deny a claim and say an idea is absurd or debatable.
\end{chapteropening}

\section[douter]{douter — to doubt}

\label{lesson:FR-C338-douter}



\begin{quote}
Before the new one: say the French for to give in, then the French for to concede.
\end{quote}

\subsection*{You'll want to know: douter}
\textbf{douter} — \textquotedblleft{}to doubt\textquotedblright{}.

\begin{quote}
Je doute qu'il vienne demain.
\end{quote}

\begin{quote}
\emph{I doubt that he will come tomorrow.}
\end{quote}

\emph{Douter que} takes the subjunctive: \emph{je doute qu'il vienne}. \emph{Se douter de} means the opposite: to suspect.

\begin{cousinweb}
\emph{douter} is inherited, through Old French, from Latin \emph{dubitāre} \textquotedblleft{}to waver, to be uncertain\textquotedblright{} (Lewis \& Short). Lewis \& Short link the Latin verb to \emph{duo} \textquotedblleft{}two\textquotedblright{}: a mind going two ways.
\end{cousinweb}

\begin{grammarlens}[title={What you've built}]
One more word for telling, explaining and describing at B1.
\end{grammarlens}

\subsection*{Your turn}
\begin{itemize}
\raggedright
  \item \emph{Say it:} \emph{douter}
  \item \emph{Say it:} \emph{douter}, once more
  \item \emph{Say it:} a sentence of your own with \emph{douter}
  \item \emph{Recall:} say the French for to give in, then the French for to concede, then say \emph{douter} again
\end{itemize}

\begin{tcolorbox}[breakable,colback=teal!4,colframe=teal!35!black,title={Before you move on}]
What does douter mean? (\textquotedblleft{}To doubt\textquotedblright{}.) Say it once more.
\end{tcolorbox}
\section[sceptique]{sceptique — sceptical}

\label{lesson:FR-C338-sceptique}



\begin{quote}
Before the new one: say the French for to concede, then the French for to doubt.
\end{quote}

\subsection*{You'll want to know: sceptique}
\textbf{sceptique} — \textquotedblleft{}sceptical\textquotedblright{}.

\begin{quote}
Mon frère reste sceptique : il ne croit pas les journaux.
\end{quote}

\begin{quote}
\emph{My brother remains sceptical: he does not believe the newspapers.}
\end{quote}

One form for both genders. The \emph{c} is silent: say \emph{septique}.

\begin{cousinweb}
\emph{sceptique} is a learned word, usually traced through Latin \emph{Sceptici} \textquotedblleft{}the Sceptics\textquotedblright{} (Lewis \& Short) to Greek \emph{skeptikós} \textquotedblleft{}inquiring, thoughtful\textquotedblright{}, which is built on \emph{sképtesthai} \textquotedblleft{}to look about, to consider\textquotedblright{}.
\end{cousinweb}

\begin{grammarlens}[title={What you've built}]
One more word for telling, explaining and describing at B1.
\end{grammarlens}

\subsection*{Your turn}
\begin{itemize}
\raggedright
  \item \emph{Say it:} \emph{sceptique}
  \item \emph{Say it:} \emph{sceptique}, once more
  \item \emph{Say it:} a sentence of your own with \emph{sceptique}
  \item \emph{Recall:} say the French for to concede, then the French for to doubt, then say \emph{sceptique} again
\end{itemize}

\begin{tcolorbox}[breakable,colback=teal!4,colframe=teal!35!black,title={Before you move on}]
What does sceptique mean? (\textquotedblleft{}Sceptical\textquotedblright{}.) Say it once more.
\end{tcolorbox}
\section[nier]{nier — to deny}

\label{lesson:FR-C338-nier}



\begin{quote}
Before the new one: say the French for to doubt, then the French for sceptical.
\end{quote}

\subsection*{You'll want to know: nier}
\textbf{nier} — \textquotedblleft{}to deny\textquotedblright{}.

\begin{quote}
Il nie avoir pris l'argent.
\end{quote}

\begin{quote}
\emph{He denies taking the money.}
\end{quote}

\emph{Nier} + infinitive, or \emph{nier que}. \emph{Je nie, nous nions}.

\begin{cousinweb}
\emph{nier} is inherited from Latin \emph{negāre} \textquotedblleft{}to say no, to deny\textquotedblright{} (Lewis \& Short), through Old French \emph{noier} (EtymDB).
\end{cousinweb}

\begin{grammarlens}[title={What you've built}]
One more word for telling, explaining and describing at B1.
\end{grammarlens}

\subsection*{Your turn}
\begin{itemize}
\raggedright
  \item \emph{Say it:} \emph{nier}
  \item \emph{Say it:} \emph{nier}, once more
  \item \emph{Say it:} a sentence of your own with \emph{nier}
  \item \emph{Recall:} say the French for to doubt, then the French for sceptical, then say \emph{nier} again
\end{itemize}

\begin{tcolorbox}[breakable,colback=teal!4,colframe=teal!35!black,title={Before you move on}]
What does nier mean? (\textquotedblleft{}To deny\textquotedblright{}.) Say it once more.
\end{tcolorbox}
\section[absurde]{absurde — absurd}

\label{lesson:FR-C338-absurde}



\begin{quote}
Before the new one: say the French for sceptical, then the French for to deny.
\end{quote}

\subsection*{You'll want to know: absurde}
\textbf{absurde} — \textquotedblleft{}absurd\textquotedblright{}.

\begin{quote}
C'est absurde de payer deux fois pour le même billet.
\end{quote}

\begin{quote}
\emph{It is absurd to pay twice for the same ticket.}
\end{quote}

One form for both genders: \emph{une idée absurde}.

\begin{cousinweb}
\emph{absurde} is a learned borrowing from Latin \emph{absurdus} \textquotedblleft{}out of tune, harsh; senseless\textquotedblright{} (EtymDB; Lewis \& Short), built on \emph{surdus} \textquotedblleft{}deaf, dull-sounding\textquotedblright{}.
\end{cousinweb}

\begin{grammarlens}[title={What you've built}]
One more word for telling, explaining and describing at B1.
\end{grammarlens}

\subsection*{Your turn}
\begin{itemize}
\raggedright
  \item \emph{Say it:} \emph{absurde}
  \item \emph{Say it:} \emph{absurde}, once more
  \item \emph{Say it:} a sentence of your own with \emph{absurde}
  \item \emph{Recall:} say the French for sceptical, then the French for to deny, then say \emph{absurde} again
\end{itemize}

\begin{tcolorbox}[breakable,colback=teal!4,colframe=teal!35!black,title={Before you move on}]
What does absurde mean? (\textquotedblleft{}Absurd\textquotedblright{}.) Say it once more.
\end{tcolorbox}
\section[discutable]{discutable — debatable, questionable}

\label{lesson:FR-C338-discutable}



\begin{quote}
Before the new one: say the French for to deny, then the French for absurd.
\end{quote}

\subsection*{You'll want to know: discutable}
\textbf{discutable} — \textquotedblleft{}debatable, questionable\textquotedblright{}.

\begin{quote}
Ta décision est discutable, mais c'est ton choix.
\end{quote}

\begin{quote}
\emph{Your decision is debatable, but it is your choice.}
\end{quote}

From \emph{discuter} + \emph{-able}: it can be argued about. \emph{C'est discutable} is a polite way to disagree.

\begin{cousinweb}
\emph{discutable} is built from \emph{discuter} and the suffix \emph{-able} (EtymDB). \emph{discuter} comes from Latin \emph{discutere} \textquotedblleft{}to strike apart, to shatter\textquotedblright{} (Lewis \& Short), built on \emph{quatere} \textquotedblleft{}to shake\textquotedblright{}.
\end{cousinweb}

\begin{grammarlens}[title={What you've built}]
One more word for telling, explaining and describing at B1.
\end{grammarlens}

\subsection*{Your turn}
\begin{itemize}
\raggedright
  \item \emph{Say it:} \emph{discutable}
  \item \emph{Say it:} \emph{discutable}, once more
  \item \emph{Say it:} a sentence of your own with \emph{discutable}
  \item \emph{Recall:} say the French for to deny, then the French for absurd, then say \emph{discutable} again
\end{itemize}

\begin{tcolorbox}[breakable,colback=teal!4,colframe=teal!35!black,title={Before you move on}]
What does discutable mean? (\textquotedblleft{}Debatable, questionable\textquotedblright{}.) Say it once more.
\end{tcolorbox}
